\documentclass[aspectratio=169,10pt]{beamer}
\usetheme[numbering=fraction,progressbar=frametitle]{metropolis}
\usepackage{fontspec}
\setsansfont{Latin Modern Sans}
\usepackage{amsmath,amssymb,mathtools,graphicx,booktabs}
\usepackage{polyglossia}
\setdefaultlanguage{french}
\setbeamertemplate{navigation symbols}{}
\setbeamerfont{frametitle}{size=\large}
\setbeamerfont{normal text}{size=\normalsize}
\setlength{\parskip}{.45em}
\hypersetup{pdfauthor={},pdfsubject={Lecture slides},pdfkeywords={},pdftitle={Variables aléatoires discrètes et loi binomiale}}
\providecommand{\N}{\mathbb{N}}
\providecommand{\Z}{\mathbb{Z}}
\providecommand{\D}{\mathbb{D}}
\providecommand{\Q}{\mathbb{Q}}
\providecommand{\R}{\mathbb{R}}
\providecommand{\C}{\mathbb{C}}
\providecommand{\vect}[1]{\overrightarrow{#1}}
\providecommand{\vi}{\vec{\imath}}
\providecommand{\vj}{\vec{\jmath}}
\providecommand{\vk}{\vec{k}}
\providecommand{\norm}[1]{\left\lVert #1\right\rVert}
\providecommand{\vc}[2]{\begin{pmatrix}#1\\#2\end{pmatrix}}
\providecommand{\deter}[2]{\det\left(#1,#2\right)}
\providecommand{\abs}[1]{\left\lvert #1\right\rvert}
\providecommand{\dd}{\,\mathrm{d}}
\providecommand{\e}{\mathrm{e}}
\providecommand{\Card}{\operatorname{Card}}
\providecommand{\E}{\operatorname{E}}
\providecommand{\V}{\operatorname{V}}
\providecommand{\degre}{^{\circ}}
\providecommand{\un}[1]{\,\mathrm{#1}}
\providecommand{\sci}[2]{#1\times10^{#2}}
\providecommand{\pgcd}{\operatorname{PGCD}}
\providecommand{\ppcm}{\operatorname{PPCM}}
\title{Variables aléatoires discrètes et loi binomiale}
\subtitle{Cours et exemples}
\date{}
\author{}
\begin{document}
\begin{frame}[plain]
\vfill
{\usebeamerfont{title}\inserttitle\par}
\medskip
{\usebeamerfont{subtitle}\insertsubtitle\par}
\medskip
{\usebeamercolor[fg]{progress bar}\rule{\textwidth}{.5pt}}
\vfill
\end{frame}
\begin{frame}{Cours de référence}
L'espérance est \(E(X)=\sum_i x_ip_i\). Elle représente la valeur moyenne à long terme lorsque l'expérience est répétée dans les mêmes conditions ; elle n'est pas nécessairement une valeur effectivement prise par \(X\).
\end{frame}
\begin{frame}{Cours de référence}
\(X\) suit une loi binomiale \(\mathcal B(n;p)\) lorsque \(X\) compte le nombre de succès dans \(n\) épreuves de Bernoulli indépendantes, de même probabilité de succès \(p\). Alors
\[
P(X=k)=\binom nk p^k(1-p)^{n-k},\qquad E(X)=np.
\]
\end{frame}
\begin{frame}{Espérance d'une variable aléatoire}
\begin{itemize}
\item Définir une loi.
\item Calculer une espérance.
\item Interpréter une moyenne théorique.

\end{itemize}
\end{frame}
\begin{frame}{Espérance d'une variable aléatoire}
Un jeu coûte \(2\) euros. Le gain brut vaut \(0\), \(3\) ou \(10\) euros avec probabilités \(0{,}6\), \(0{,}3\), \(0{,}1\). Construire la loi du gain net \(X\), calculer \(E(X)\) et décider si le jeu est favorable au joueur.
\end{frame}
\begin{frame}{Espérance d'une variable aléatoire}
La loi de \(X\) associe à chaque valeur \(x_i\) la probabilité \(p_i\). L'espérance \(E(X)=\sum x_ip_i\) est une moyenne théorique à long terme.
\end{frame}
\begin{frame}{Espérance d'une variable aléatoire}
Reprendre les données de l'activité et rédiger la solution dans l'ordre : définition des objets, propriété utilisée, calcul exact ou approché, puis phrase de réponse. La réponse modèle figure dans l'édition enseignant.
\end{frame}
\begin{frame}{Espérance d'une variable aléatoire}
\textbf{Conduite en 1 h :} 7 min d'automatismes ; 10 min de mise en commun et trace ; 28 min d'exercices essentiels ; 5 min de correction ; 3 min de billet de sortie.

\textbf{Conduite en 2 h :} ajouter 20 min de recherche de l'activité, 20 min d'institutionnalisation, un second exercice autonome et 10 min de correction comparée.

\textbf{Différenciation :} fournir aux élèves fragiles une fiche à quatre cases (données, modèle, calcul, conclusion). Demander aux élèves rapides de produire un contre-exemple, une justification ou une seconde méthode.
\end{frame}
\begin{frame}{Loi binomiale : reconnaître une situation}
\begin{itemize}
\item Reconnaître un schéma de Bernoulli.
\item Identifier \(n\) et \(p\).
\item Définir \(X\).

\end{itemize}
\end{frame}
\begin{frame}{Loi binomiale : reconnaître une situation}
Un client accepte une offre avec probabilité \(0{,}25\). Un conseiller contacte indépendamment huit clients. Définir \(X\), justifier la loi binomiale et identifier \(n\) et \(p\).
\end{frame}
\begin{frame}{Loi binomiale : reconnaître une situation}
\(X\sim\mathcal B(n;p)\) si \(X\) compte les succès dans \(n\) épreuves indépendantes de même probabilité \(p\).
\end{frame}
\begin{frame}{Loi binomiale : reconnaître une situation}
Reprendre les données de l'activité et rédiger la solution dans l'ordre : définition des objets, propriété utilisée, calcul exact ou approché, puis phrase de réponse. La réponse modèle figure dans l'édition enseignant.
\end{frame}
\begin{frame}{Loi binomiale : reconnaître une situation}
\textbf{Conduite en 1 h :} 7 min d'automatismes ; 10 min de mise en commun et trace ; 28 min d'exercices essentiels ; 5 min de correction ; 3 min de billet de sortie.

\textbf{Conduite en 2 h :} ajouter 20 min de recherche de l'activité, 20 min d'institutionnalisation, un second exercice autonome et 10 min de correction comparée.

\textbf{Différenciation :} fournir aux élèves fragiles une fiche à quatre cases (données, modèle, calcul, conclusion). Demander aux élèves rapides de produire un contre-exemple, une justification ou une seconde méthode.
\end{frame}
\begin{frame}{Triangle de Pascal et coefficients binomiaux}
\begin{itemize}
\item Construire le triangle de Pascal.
\item Lire \(\binom nk\).
\item Interpréter un coefficient.

\end{itemize}
\end{frame}
\begin{frame}{Triangle de Pascal et coefficients binomiaux}
Construire les lignes \(n=0\) à \(n=6\) du triangle de Pascal. Expliquer le calcul du nombre \(15\) de la ligne \(6\), puis lire \(\binom62\) et \(\binom64\).
centertabular*9c\&\&\&\&1\&\&\&\&

\&\&\&1\&\&1\&\&\&

\&\&1\&\&2\&\&1\&\&

\&1\&\&3\&\&3\&\&1\&

1\&\&4\&\&6\&\&4\&\&1tabularcenter
\end{frame}
\begin{frame}{Triangle de Pascal et coefficients binomiaux}
Le triangle de Pascal vérifie \(\binom nk=\binom{n-1}{k-1}+\binom{n-1}{k}\). Les bords valent \(1\).
\end{frame}
\begin{frame}{Triangle de Pascal et coefficients binomiaux}
Reprendre les données de l'activité et rédiger la solution dans l'ordre : définition des objets, propriété utilisée, calcul exact ou approché, puis phrase de réponse. La réponse modèle figure dans l'édition enseignant.
\end{frame}
\begin{frame}{Triangle de Pascal et coefficients binomiaux}
\textbf{Conduite en 1 h :} 7 min d'automatismes ; 10 min de mise en commun et trace ; 28 min d'exercices essentiels ; 5 min de correction ; 3 min de billet de sortie.

\textbf{Conduite en 2 h :} ajouter 20 min de recherche de l'activité, 20 min d'institutionnalisation, un second exercice autonome et 10 min de correction comparée.

\textbf{Différenciation :} fournir aux élèves fragiles une fiche à quatre cases (données, modèle, calcul, conclusion). Demander aux élèves rapides de produire un contre-exemple, une justification ou une seconde méthode.
\end{frame}
\begin{frame}{Probabilités binomiales}
\begin{itemize}
\item Calculer \(P(X=k)\).
\item Traduire au moins/au plus.
\item Utiliser le complémentaire.

\end{itemize}
\end{frame}
\begin{frame}{Probabilités binomiales}
\begin{columns}[T,onlytextwidth]
\begin{column}{.57\textwidth}
Pour \(X\sim\mathcal B(8;0{,}25)\), calculer \(P(X=2)\), \(P(X=0)\) puis \(P(X\geq1)\). Traduire chaque événement en phrase.
center center
\end{column}
\begin{column}{.4\textwidth}
\centering
\includegraphics[width=\linewidth,height=4.5cm,keepaspectratio]{figures/stmg-07-3.pdf}
\end{column}
\end{columns}
\end{frame}
\begin{frame}{Probabilités binomiales}
Si \(X\sim\mathcal B(n;p)\), alors \(P(X=k)=\binom nkp^k(1-p)^{n-k}\). Traduire soigneusement « au moins », « au plus » et « aucun ».
\end{frame}
\begin{frame}{Probabilités binomiales}
Reprendre les données de l'activité et rédiger la solution dans l'ordre : définition des objets, propriété utilisée, calcul exact ou approché, puis phrase de réponse. La réponse modèle figure dans l'édition enseignant.
\end{frame}
\begin{frame}{Probabilités binomiales}
\textbf{Conduite en 1 h :} 7 min d'automatismes ; 10 min de mise en commun et trace ; 28 min d'exercices essentiels ; 5 min de correction ; 3 min de billet de sortie.

\textbf{Conduite en 2 h :} ajouter 20 min de recherche de l'activité, 20 min d'institutionnalisation, un second exercice autonome et 10 min de correction comparée.

\textbf{Différenciation :} fournir aux élèves fragiles une fiche à quatre cases (données, modèle, calcul, conclusion). Demander aux élèves rapides de produire un contre-exemple, une justification ou une seconde méthode.
\end{frame}
\begin{frame}{Bilan du chapitre et évaluation}
\begin{itemize}
\item Calculer et interpréter \(E(X)\).
\item Résoudre un problème binomial.
\item Analyser ses erreurs.

\end{itemize}
\end{frame}
\begin{frame}{Bilan du chapitre et évaluation}
Un centre reçoit \(120\) groupes de \(10\) clients, chaque client réclamant avec probabilité \(0{,}08\). Estimer le nombre moyen de réclamations par groupe, puis sur l'ensemble des groupes. Discuter ce que l'espérance permet ou non de prévoir.
\end{frame}
\begin{frame}{Bilan du chapitre et évaluation}
Pour \(X\sim\mathcal B(n;p)\), \(E(X)=np\). La valeur donne le nombre moyen de succès sur un grand nombre de séries de \(n\) épreuves.
\end{frame}
\begin{frame}{Bilan du chapitre et évaluation}
Reprendre les données de l'activité et rédiger la solution dans l'ordre : définition des objets, propriété utilisée, calcul exact ou approché, puis phrase de réponse. La réponse modèle figure dans l'édition enseignant.
\end{frame}
\begin{frame}{Bilan du chapitre et évaluation}
\textbf{Conduite en 1 h :} 7 min d'automatismes ; 10 min de mise en commun et trace ; 28 min d'exercices essentiels ; 5 min de correction ; 3 min de billet de sortie.

\textbf{Conduite en 2 h :} ajouter 20 min de recherche de l'activité, 20 min d'institutionnalisation, un second exercice autonome et 10 min de correction comparée.

\textbf{Différenciation :} fournir aux élèves fragiles une fiche à quatre cases (données, modèle, calcul, conclusion). Demander aux élèves rapides de produire un contre-exemple, une justification ou une seconde méthode.
\end{frame}
\end{document}
